% figures/w04-los-geometry.svg — LOS 유도의 기하. Fossen 의 표준 그림 배치.
%
%   matlab -batch "addpath _tools; w04_losgeo"          <- 먼저 기하를 계산하고
%   bash _tools/tikz2svg.sh figures/src/w04-los-geometry.tex figures/w04-los-geometry.svg
%
% 배치와 기호는 Fossen 의 TTK 4190 "Guidance, Navigation and Control of
% Vehicles" 강의노트, 그리고 Handbook of Marine Craft Hydrodynamics and Motion
% Control (2021, 2nd ed.) §12.3 을 따른다. 사용자가 그 그림으로 맞춰 달라고
% 지시했다 (2026-09-08).
%
%   p_i^n, p_{i+1}^n   웨이포인트 (NED). 굵은 소문자에 위첨자 n
%   pi_p               경로 접선각
%   x_e^p, y_e^p       경로좌표계 {p} 에서의 along-track / cross-track 오차
%   Delta              look-ahead distance. 수선의 발에서 **경로를 따라** 잰다
%   chi_d              요구 침로각 = pi_p - atan(y_e^p/Delta)
%   chi, psi, beta_c   실제 침로각, 선수각, 크랩각.  chi = psi + beta_c
%
% 좌표는 전부 _tools/w04_losgeo.m 이 계산한다. **이 파일에는 산술이 없다** —
% TikZ 의 calc 에 매크로 좌표와 음수 계수를 함께 물리면 pgf 의 계수 파서가
% 폭주하고, 이름표를 눈으로 놓으면 세 개가 LOS 벡터 위에 겹친다. 둘 다 실제로
% 겪었다. 설명 글은 그림 **아래**에 둔다 — 옆에 두면 전체 폭이 A4 본문 폭
% 703 px 를 넘어 인쇄될 때 축소된다 (CLAUDE.md §4 규칙 6-0).
\documentclass[border=5pt]{standalone}
% gnc-style.tex — 이 볼트의 TikZ 그림이 공유하는 스타일.
%
% 저널 그림 규약을 스타일 하나로 굳혀 둔다. 그림마다 선 굵기나 화살촉을
% 다시 정하지 않는다 — 그것이 그림이 제각각으로 보이는 원인이다.
%
% 쓰는 법:  % gnc-style.tex — 이 볼트의 TikZ 그림이 공유하는 스타일.
%
% 저널 그림 규약을 스타일 하나로 굳혀 둔다. 그림마다 선 굵기나 화살촉을
% 다시 정하지 않는다 — 그것이 그림이 제각각으로 보이는 원인이다.
%
% 쓰는 법:  % gnc-style.tex — 이 볼트의 TikZ 그림이 공유하는 스타일.
%
% 저널 그림 규약을 스타일 하나로 굳혀 둔다. 그림마다 선 굵기나 화살촉을
% 다시 정하지 않는다 — 그것이 그림이 제각각으로 보이는 원인이다.
%
% 쓰는 법:  \input{gnc-style}  (figures/src/ 안에서)

\usepackage{tikz}
\usepackage{amsmath}
\usepackage{xcolor}
\usetikzlibrary{arrows.meta, positioning, calc, fit, backgrounds, decorations.pathmorphing}

% ---- 색 --------------------------------------------------------------------
% 색은 강조 하나에만 쓴다. 나머지는 검정.
\definecolor{ink}{HTML}{1A1A1A}
\definecolor{accent}{HTML}{7C3AED}
\definecolor{warn}{HTML}{B91C1C}
\definecolor{soft}{HTML}{666666}

% ---- 선과 화살촉 -----------------------------------------------------------
\tikzset{
  % 신호선. 화살촉은 작고 뾰족한 것 하나뿐이다.
  sig/.style   = {ink, line width=0.5pt, -{Stealth[length=4pt, width=3pt]}},
  wire/.style  = {ink, line width=0.5pt},
  % 강조 경로 — 파선으로 구분한다. 흑백 인쇄에서도 살아야 하기 때문이다.
  asig/.style  = {accent, line width=0.5pt, densely dashed,
                  -{Stealth[length=4pt, width=3pt]}},
  awire/.style = {accent, line width=0.5pt, densely dashed},
  % 블록. 흰 바탕, 직각, 검은 테두리.
  blk/.style   = {draw=ink, line width=0.55pt, fill=white,
                  minimum width=17mm, minimum height=8mm, inner sep=2pt, font=\small},
  % 합산점
  sum/.style   = {draw=ink, line width=0.55pt, fill=white, circle,
                  inner sep=0pt, minimum size=4.6mm, font=\scriptsize},
  asum/.style  = {draw=accent, line width=0.55pt, fill=white, circle,
                  inner sep=0pt, minimum size=4.6mm, font=\scriptsize},
  % 분기점
  dot/.style   = {ink, fill=ink, circle, inner sep=0pt, minimum size=1.5mm},
  % 신호 이름 — 선 위에 작게
  sname/.style = {font=\small, inner sep=1.5pt},
  % 그림 안의 짧은 설명. 그림 안의 글은 최소로 둔다
  note/.style  = {font=\scriptsize, soft, inner sep=1.5pt},
  anote/.style = {font=\scriptsize, accent, inner sep=1.5pt},
  % 축
  axis/.style  = {ink, line width=0.5pt},
  tick/.style  = {font=\scriptsize, soft},
}

% 캡션 한 줄 — 그림 아래 가는 선과 함께
\newcommand{\gnccaption}[2]{%
  \node[anchor=north west, text width=#1, align=left, font=\scriptsize, ink]
        at (current bounding box.south west) {\vspace{2pt}#2};}
  (figures/src/ 안에서)

\usepackage{tikz}
\usepackage{amsmath}
\usepackage{xcolor}
\usetikzlibrary{arrows.meta, positioning, calc, fit, backgrounds, decorations.pathmorphing}

% ---- 색 --------------------------------------------------------------------
% 색은 강조 하나에만 쓴다. 나머지는 검정.
\definecolor{ink}{HTML}{1A1A1A}
\definecolor{accent}{HTML}{7C3AED}
\definecolor{warn}{HTML}{B91C1C}
\definecolor{soft}{HTML}{666666}

% ---- 선과 화살촉 -----------------------------------------------------------
\tikzset{
  % 신호선. 화살촉은 작고 뾰족한 것 하나뿐이다.
  sig/.style   = {ink, line width=0.5pt, -{Stealth[length=4pt, width=3pt]}},
  wire/.style  = {ink, line width=0.5pt},
  % 강조 경로 — 파선으로 구분한다. 흑백 인쇄에서도 살아야 하기 때문이다.
  asig/.style  = {accent, line width=0.5pt, densely dashed,
                  -{Stealth[length=4pt, width=3pt]}},
  awire/.style = {accent, line width=0.5pt, densely dashed},
  % 블록. 흰 바탕, 직각, 검은 테두리.
  blk/.style   = {draw=ink, line width=0.55pt, fill=white,
                  minimum width=17mm, minimum height=8mm, inner sep=2pt, font=\small},
  % 합산점
  sum/.style   = {draw=ink, line width=0.55pt, fill=white, circle,
                  inner sep=0pt, minimum size=4.6mm, font=\scriptsize},
  asum/.style  = {draw=accent, line width=0.55pt, fill=white, circle,
                  inner sep=0pt, minimum size=4.6mm, font=\scriptsize},
  % 분기점
  dot/.style   = {ink, fill=ink, circle, inner sep=0pt, minimum size=1.5mm},
  % 신호 이름 — 선 위에 작게
  sname/.style = {font=\small, inner sep=1.5pt},
  % 그림 안의 짧은 설명. 그림 안의 글은 최소로 둔다
  note/.style  = {font=\scriptsize, soft, inner sep=1.5pt},
  anote/.style = {font=\scriptsize, accent, inner sep=1.5pt},
  % 축
  axis/.style  = {ink, line width=0.5pt},
  tick/.style  = {font=\scriptsize, soft},
}

% 캡션 한 줄 — 그림 아래 가는 선과 함께
\newcommand{\gnccaption}[2]{%
  \node[anchor=north west, text width=#1, align=left, font=\scriptsize, ink]
        at (current bounding box.south west) {\vspace{2pt}#2};}
  (figures/src/ 안에서)

\usepackage{tikz}
\usepackage{amsmath}
\usepackage{xcolor}
\usetikzlibrary{arrows.meta, positioning, calc, fit, backgrounds, decorations.pathmorphing}

% ---- 색 --------------------------------------------------------------------
% 색은 강조 하나에만 쓴다. 나머지는 검정.
\definecolor{ink}{HTML}{1A1A1A}
\definecolor{accent}{HTML}{7C3AED}
\definecolor{warn}{HTML}{B91C1C}
\definecolor{soft}{HTML}{666666}

% ---- 선과 화살촉 -----------------------------------------------------------
\tikzset{
  % 신호선. 화살촉은 작고 뾰족한 것 하나뿐이다.
  sig/.style   = {ink, line width=0.5pt, -{Stealth[length=4pt, width=3pt]}},
  wire/.style  = {ink, line width=0.5pt},
  % 강조 경로 — 파선으로 구분한다. 흑백 인쇄에서도 살아야 하기 때문이다.
  asig/.style  = {accent, line width=0.5pt, densely dashed,
                  -{Stealth[length=4pt, width=3pt]}},
  awire/.style = {accent, line width=0.5pt, densely dashed},
  % 블록. 흰 바탕, 직각, 검은 테두리.
  blk/.style   = {draw=ink, line width=0.55pt, fill=white,
                  minimum width=17mm, minimum height=8mm, inner sep=2pt, font=\small},
  % 합산점
  sum/.style   = {draw=ink, line width=0.55pt, fill=white, circle,
                  inner sep=0pt, minimum size=4.6mm, font=\scriptsize},
  asum/.style  = {draw=accent, line width=0.55pt, fill=white, circle,
                  inner sep=0pt, minimum size=4.6mm, font=\scriptsize},
  % 분기점
  dot/.style   = {ink, fill=ink, circle, inner sep=0pt, minimum size=1.5mm},
  % 신호 이름 — 선 위에 작게
  sname/.style = {font=\small, inner sep=1.5pt},
  % 그림 안의 짧은 설명. 그림 안의 글은 최소로 둔다
  note/.style  = {font=\scriptsize, soft, inner sep=1.5pt},
  anote/.style = {font=\scriptsize, accent, inner sep=1.5pt},
  % 축
  axis/.style  = {ink, line width=0.5pt},
  tick/.style  = {font=\scriptsize, soft},
}

% 캡션 한 줄 — 그림 아래 가는 선과 함께
\newcommand{\gnccaption}[2]{%
  \node[anchor=north west, text width=#1, align=left, font=\scriptsize, ink]
        at (current bounding box.south west) {\vspace{2pt}#2};}

% figures/src/w04-los-geometry-geo.tex — 손으로 고치지 않는다.
%
%   matlab -batch "addpath _tools; w04_losgeo"
%
% Fossen TTK4190 의 표준 LOS 그림 배치. 좌표는 (East, North) 순이고
% 단위는 mm (1 m = 1.55 mm).

\def\WPi{(0.000,0.000)}
\def\WPnext{(118.737,99.632)}
\def\Foot{(37.996,31.882)}
\def\Ves{(53.937,12.884)}
\def\Aim{(83.116,69.742)}
\def\Cross{(53.937,45.258)}
\def\PathA{(-9.499,-7.971)}
\def\PathB{(132.985,111.588)}
\def\MidXe{(18.998,15.941)}
\def\MidYe{(45.966,22.383)}
\def\MidDel{(60.556,50.812)}
\def\MidLOS{(69.110,42.451)}
\def\RaA{(36.004,30.211)}
\def\RaB{(37.675,28.219)}
\def\RaC{(39.667,29.891)}
\def\DropA{(83.116,-13.116)}
\def\DropB{(53.937,-13.116)}
\def\RbA{(80.116,-13.116)}
\def\RbB{(80.116,-10.116)}
\def\RbC{(83.116,-10.116)}
\def\Xntop{(53.937,58.884)}
\def\RadPsi{20.00}
\def\ArcTopPsi{(53.937,32.884)}
\def\LabPsi{(56.886,36.904)}
\def\RadChi{26.00}
\def\ArcTopChi{(53.937,38.884)}
\def\LabChi{(59.699,42.530)}
\def\RadChid{33.00}
\def\ArcTopChid{(53.937,45.884)}
\def\LabChid{(62.674,49.044)}
\def\ArcTopBeta{(57.324,26.469)}
\def\LabBeta{(59.654,30.479)}
\def\ArcTopCorr{(76.758,64.407)}
\def\RadCorr{8.30}
\def\ArcMidCorr{(76.942,62.004)}
\def\HullBow{(55.388,18.706)}
\def\HullStbd{(55.559,9.883)}
\def\HullPort{(51.096,10.996)}
\def\LabXb{(60.590,39.567)}
\def\LabU{(69.296,50.899)}
\def\LabDelta{(67.525,62.665)}
\def\LabPip{(49.937,48.258)}

\def\ScrPip{40.0000}
\def\ScrChid{62.8337}
\def\ScrChi{68.0000}
\def\ScrPsi{76.0000}

\def\ValPip{50}
\def\ValChid{27.17}
\def\ValChi{22}
\def\ValPsi{14}
\def\ValBeta{8}
\def\ValXe{32}
\def\ValYe{16}
\def\ValDelta{38}
\def\ValCorr{22.83}
\def\ValLOS{41.23}
\def\Scale{1.55}


%  색은 \tikzset 보다 먼저 정의해야 한다 — 스타일이 곧바로 쓴다.
\definecolor{knownc}{HTML}{1D4ED8}

\tikzset{
  pathline/.style = {ink, line width=0.9pt},
  aux/.style      = {soft, line width=0.35pt, densely dashed},
  errline/.style  = {errc, line width=1.0pt},
  alngline/.style = {alng, line width=1.0pt},
  losline/.style  = {accent, line width=1.0pt, -{Stealth[length=4.6pt,width=3.4pt]}},
  velline/.style  = {pathc, line width=1.0pt, -{Stealth[length=4.6pt,width=3.4pt]}},
  bodyline/.style = {ink, line width=0.8pt, -{Stealth[length=4.2pt,width=3.0pt]}},
  ang/.style      = {soft, line width=0.5pt},
  rmark/.style    = {soft, line width=0.4pt},
  %  강의 슬라이드처럼 **식을 그림 안에 크게** 둔다. 상자에 넣어야 그림의
  %  선들과 섞이지 않는다. 색은 그 식이 무엇인지에 따라 다르다:
  %    known  이미 아는 값        knownc 파랑
  %    want   구하려는 값         accent 보라
  %    limit  극한에서 되는 것    pathc  청록
  known/.style = {draw=knownc, line width=0.6pt, fill=white, inner sep=3.2pt,
                  rounded corners=1pt, font=\small},
  want/.style  = {draw=accent, line width=0.9pt, fill=white, inner sep=3.4pt,
                  rounded corners=1pt, font=\small},
  limit/.style = {draw=pathc, line width=0.6pt, fill=white, inner sep=3.2pt,
                  rounded corners=1pt, font=\small},
  capt/.style  = {font=\scriptsize\bfseries, warn},
  lead/.style  = {soft, line width=0.35pt, -{Stealth[length=3.4pt,width=2.4pt]}},
}

\begin{document}
\begin{tikzpicture}[x=1mm, y=1mm]

% ---- 북·동 축 ------------------------------------------------------------
\draw[axis, -{Stealth[length=4pt,width=3pt]}] (2,96) -- (2,109);
\node[font=\scriptsize, anchor=south] at (2,109.5) {North};
\draw[axis, -{Stealth[length=4pt,width=3pt]}] (2,96) -- (15,96);
\node[font=\scriptsize, anchor=west] at (15.5,96) {East};

% ---- 경로 ----------------------------------------------------------------
\draw[pathline] \PathA -- \PathB;
\fill[ink] \WPi    circle (1.5pt);
\fill[ink] \WPnext circle (1.5pt);
\node[font=\small, anchor=east]        at (-3.0,1.0)    {$\mathbf{p}_{i}^{n}$};
%  p_i 의 오른쪽 아래에 둔다. 왼쪽에 두면 그림 왼쪽 끝에 딱 붙어 잘려 보인다.
\node[note, anchor=north west]         at (2.5,-2.0)    {$[\,x_i^n,\ y_i^n\,]^{\!\top}$};
\node[font=\small, anchor=south east]  at (117.0,104.5) {$\mathbf{p}_{i+1}^{n}$};
%  두 웨이포인트 모두 성분을 적는다. 하나만 적으면 다른 하나가 무엇인지
%  독자가 유추해야 한다 — 사용자 지적 2026-09-08.
%  경로선 **아래** 오른쪽에 둔다. 선 위에 두면 글자가 선에 물린다.
\node[note, anchor=north west]         at (119.5,96.0)  {$[\,x_{i+1}^n,\ y_{i+1}^n\,]^{\!\top}$};

% ---- 선박을 지나는 북쪽 점선, 그리고 pi_p ------------------------------
\draw[aux] \Ves -- \Xntop;
\node[font=\small, soft, anchor=south east] at (53.0,59.0) {$x_n$};
\draw[ang] (53.937,58.258) arc (90:\ScrPip:13);
\node[font=\small, anchor=east] at \LabPip {$\pi_p$};

% ---- 수선의 발, along-track 과 cross-track -----------------------------
\draw[alngline] \WPi -- \Foot;
\draw[errline]  \Foot -- \Ves;
\fill[ink] \Foot circle (1.1pt);
\draw[rmark] \RaA -- \RaB -- \RaC;

\node[alng, font=\small, anchor=south east] at (17.0,17.8) {$x_e^{p}$};
\node[errc, font=\small, anchor=east]       at (43.6,22.4) {$y_e^{p}$};

% ---- look-ahead 와 조준점 -----------------------------------------------
\draw[alngline] \Foot -- \Aim;
\fill[accent] \Aim circle (1.4pt);
\node[alng, font=\small, anchor=south east] at \LabDelta {$\Delta$};
%  Delta 가 무엇인지 말로도 짚어 준다 — 슬라이드의 "look-ahead distance".
\node[capt,anchor=east] at (44.0,74.0) {look-ahead distance};
\draw[lead] (45.0,73.4) -- (57.0,60.0);

% ---- LOS 벡터 ------------------------------------------------------------
\draw[losline] \Ves -- \Aim;
\node[accent, font=\scriptsize, anchor=north west, rotate=\ScrChid]
     at (70.8,41.0) {LOS vector};

%  조준점에서의 각. 두 선을 조준점에서 **뒤로** 본 구간이다.
\draw[ang] \ArcTopCorr arc (\ScrPip+180:\ScrChid+180:\RadCorr);
\draw[soft, line width=0.35pt, -{Stealth[length=3.4pt,width=2.4pt]}]
     (98.0,52.0) -- \ArcMidCorr;
\node[font=\scriptsize, anchor=west] at (98.5,51.6)
     {$\tan^{-1}\!\big(y_e^{p}/\Delta\big)$};

%  조준점에서 내린 세로 점선, 선박에서 오는 가로 점선, 만나는 곳의 직각
\draw[aux] \Aim   -- \DropA;
\draw[aux] \Ves   -- \DropB;
\draw[rmark] \RbA -- \RbB -- \RbC;

% ---- 선박 ----------------------------------------------------------------
\fill[ink, opacity=0.85] \HullBow -- \HullStbd -- \HullPort -- cycle;

\draw[bodyline] \Ves -- ++(\ScrPsi:24);
\node[font=\small] at \LabXb {$x_b$};

\draw[velline] \Ves -- ++(\ScrChi:38);
\node[pathc, font=\small] at \LabU {$U$};

\draw[ang] \ArcTopPsi  arc (90:\ScrPsi:\RadPsi);
\node[font=\small]       at \LabPsi  {$\psi$};
\draw[ang] \ArcTopChi  arc (90:\ScrChi:\RadChi);
\node[pathc, font=\small] at \LabChi {$\chi$};
\draw[ang] \ArcTopChid arc (90:\ScrChid:\RadChid);
\node[accent, font=\small] at \LabChid {$\chi_d$};

\draw[accent, line width=0.7pt] \ArcTopBeta arc (\ScrPsi:\ScrChi:14);
\node[anote] at \LabBeta {$\beta_c$};

\node[note, anchor=north east] at (51.0,8.0) {$(x^n,\ y^n)$};

% ===================== 그림 안의 식 ======================================
%  강의 슬라이드의 배치를 따른다: 아는 값은 파란 상자, 구하려는 값은 보라
%  상자, 극한에서 성립하는 것은 청록 상자. 상자마다 무엇인지 한 낱말로 적는다.

%  경로 접선각 — 두 웨이포인트에서 바로 나온다
\node[known, anchor=north west] at (-6,-32)
     {$\pi_p = \operatorname{atan2}\!\big(y_{i+1}^n - y_i^n,\ x_{i+1}^n - x_i^n\big)$};
\node[capt,anchor=north west] at (-6,-40.5) {Known};

%  cross-track 오차 — 위치와 경로에서 바로 나온다
\node[known, anchor=north west] at (-6,-48)
     {$y_e^{p} = -\big(x^n - x_i^n\big)\sin\pi_p + \big(y^n - y_i^n\big)\cos\pi_p$};
\node[capt,anchor=north west] at (-6,-56.5) {Known};

%  웨이포인트 둘
\node[known, anchor=north west] at (-6,-64)
     {$\mathbf{p}_i^{n} = [\,x_i^n,\ y_i^n\,]^{\!\top}$,\quad
       $\mathbf{p}_{i+1}^{n} = [\,x_{i+1}^n,\ y_{i+1}^n\,]^{\!\top}$};
\node[capt,anchor=north west] at (-6,-72.5) {Known};

%  구하려는 것
\node[want, anchor=north west] at (76,-32)
     {$\chi_d = \pi_p - \tan^{-1}\!\big(K_p\,y_e^{p}\big),
       \qquad K_p = \dfrac{1}{\Delta}$};
\node[capt,anchor=north west] at (76,-46) {Want to compute};

%  극한에서 성립하는 것
\node[limit, anchor=north west] at (76,-56)
     {$y_e^{p} \to 0, \qquad \chi_d \to \pi_p$};
\node[anchor=north west, align=left, text width=64mm, font=\scriptsize, soft] at (76,-65)
     {On the path the correction vanishes and the vessel steers straight along
      the leg. That is the whole objective, written as a limit.};

% ===================== 아래 : 법칙과 수치 ================================
\draw[soft, line width=0.3pt] (-6,-80) -- (146,-80);

\node[anchor=north west, align=left, text width=68mm, font=\scriptsize] at (-6,-84) {%
  \textbf{Two terms, and that is all.} $\pi_p$ lines the vessel up \emph{with}
  the path; $-\tan^{-1}(y_e^{p}/\Delta)$ leans it \emph{towards} the path.
  Nothing else appears, which is exactly why §4-7 can say what the law cannot
  do.\\[4pt]
  \textbf{The correction is bounded by $90^\circ$.} As $y_e^{p}\to 0$ the
  vessel is on the path and $\chi_d\to\pi_p$; as $y_e^{p}\to\infty$ the law
  gives $\chi_d\to\pi_p-90^\circ$, which heads straight at the path and never
  away from it.\\[4pt]
  \textbf{$\Delta$ is measured along the path}, from the foot of the
  perpendicular and not from the vessel. The vessel's own distance to the aim
  point is $\sqrt{\Delta^2+(y_e^{p})^2}$, which is larger.};

\node[anchor=north west, align=left, text width=68mm, font=\scriptsize] at (76,-84) {%
  {\small\bfseries Notation, and the numbers drawn here}\\[4pt]
  Waypoints are $\mathbf{p}_i^{n} = [\,x_i^n,\ y_i^n\,]^{\!\top}$, resolved in
  $\{n\}$. The two errors carry the superscript $p$ because they are resolved
  in the \textbf{path} frame $\{p\}$ — not in $\{n\}$ and not in $\{b\}$. The
  superscripts are what keep the three frames apart.\\[4pt]
  Drawn with $\pi_p = \ValPip^\circ$, $x_e^{p} = \ValXe$\,m,
  $y_e^{p} = \ValYe$\,m and $\Delta = \ValDelta$\,m, so the correction is
  $\tan^{-1}(\ValYe/\ValDelta) = \ValCorr^\circ$ and
  $\chi_d = \ValChid^\circ$; the line-of-sight distance is $\ValLOS$\,m.
  The hull is at $\psi = \ValPsi^\circ$ with $\beta_c = \ValBeta^\circ$, so
  $\chi = \ValChi^\circ$ — not yet on the commanded course, which is what
  makes the three arcs at the hull distinct.\\[4pt]
  \textbf{Course or heading.} The law as written commands a \emph{course}
  $\chi_d$. This course carries a heading autopilot from Week~3, so §4-4
  commands $\psi_d$: the same expression with $\psi_d$ in place of $\chi_d$.
  The two agree exactly when $\beta_c = 0$, and §4-7 measures what the
  difference costs when it is not.};

\end{tikzpicture}
\end{document}
